% Chapter 15: Ground-Reaction Forces and Constrained Dynamics
\chapter[Ground-Reaction Forces and Constrained Dynamics]{Ground-Reaction Forces\\and Constrained Dynamics}\label{ch:grf}

Ground-reaction force (GRF) is the force exerted by the ground on a body in contact with it. In golf, the practically useful observable is usually a \textbf{ground-reaction wrench}: three force components and three moments expressed about a declared force-plate origin. That wrench changes whole-body linear and angular momentum. It is not a direct neural command, a direct measure of muscular effort, or by itself a measure of energy delivered to the club.

\section{Scope of This Chapter}\label{grf-scope-of-this-chapter}

This chapter distinguishes three operations that are often conflated:

\begin{enumerate}
\def\labelenumi{\arabic{enumi}.}
\item
  \textbf{Measurement:} a force plate measures an external contact wrench.
\item
  \textbf{Inverse dynamics:} a model reconciles measured kinematics, inertia, and external wrenches to estimate generalized forces \citep{Nesbit2005, Gatt1998}.
\item
  \textbf{Counterfactual attribution:} the same declared model is reevaluated at a measured state after selected velocity or control terms are set to zero.
\end{enumerate}

The third operation is model dependent. It does not turn a force-plate trace into a unique estimate of muscle activation.

The central question is how internally generated motion and external support jointly shape club delivery. A foot constraint changes the accelerations that joint inputs can produce; those accelerations determine the required contact reaction. Contact and muscle action therefore cannot be assigned independent roles merely by naming one a reaction and the other a control. Momentum, work, and input attribution describe different aspects of the same motion.

\section{Declare the System and the Frame}\label{sec:grf-definition}\label{grf-declare-the-system-and-the-frame}

Use an inertial laboratory frame with \(z\) upward. Declare the horizontal axes explicitly, for example \(x\) along the target line and \(y\) completing a right-handed frame. Anatomical anterior-posterior and medial-lateral directions need a stated transformation; their names do not automatically identify the laboratory axes for a golfer whose body rotates.

Unless stated otherwise, the system below includes the \textbf{golfer and club}, ends before ball contact, and neglects aerodynamic loads. Hand-club interaction forces and couples are then internal and cancel in whole-system balances. If the club is excluded, its wrench on the hands is external and must be included. During ball impact, include the ball's contact impulse or enlarge the system again. Use the mass and center of mass of the chosen system consistently.

The third-law partner of the ground's force on the shoe is the shoe's force on the ground. Gravity and the normal reaction both act on the golfer, so they are \textbf{not} a third-law pair. Their equality in magnitude in a simple static stance follows from zero acceleration. Dynamic contact need not equal weight.

\section[Force, Moment, Center of Pressure, and Free Moment]{Force, Moment, Center of Pressure,\\and Free Moment}\label{grf-force-moment-center-of-pressure-and-free-moment}

Let a force plate report force

\[
\bm{F}=(F_x,F_y,F_z)
\]

and moment about its origin

\[
\bm{M}_O=(M_x,M_y,M_z).
\]

The axes, signs, origin, and units must accompany the data. Suppose the flat contact surface is \(z=z_s\) relative to the reporting origin \(O\). Transport the moment to \(C=(x_C,y_C,z_s)\) using

\[
\bm M_C=\bm M_O-\bm r_{OC}\times\bm F.
\]

The \textbf{center of pressure} (COP) is the surface location where the two tangential moment components vanish. For sufficiently large \(F_z\),

\[
\begin{aligned}
x_C&=\frac{z_sF_x-M_y}{F_z},\\
y_C&=\frac{M_x+z_sF_y}{F_z},\\
M_{C,z}&=M_z-(x_CF_y-y_CF_x).
\end{aligned}
\]

At a surface origin, \(z_s=0\), these reduce to \(x_C=-M_y/F_z\) and \(y_C=M_x/F_z\). The remaining \(M_{C,z}\) is the vertical \textbf{free moment}. Thus the complete wrench generally requires a force \textbf{and} a couple at the COP: not all moments vanish there. The free moment represents a distributed contact couple and cannot be reconstructed from a planar point-contact force alone. COP is also a ratio, so it becomes unstable as \(F_z\) approaches zero. Low-load samples need an explicit mask rather than a silently reported COP.

For example, a surface point \(C=(0.1,-0.05,0.04)\) m, force \(\bm F=(100,-50,800)\) N and free moment \(M_{C,z}=8\) N m produce \(\bm M_O=(-38,-76,8)\) N m. Omitting the 0.04 m surface offset gives the wrong COP even though the force and moment data are otherwise correct.

Pressure insoles measure plantar pressure, primarily the normal component. Shear forces and free moment require additional sensing or an empirically validated estimator. A study that predicts six-axis force-plate quantities from pressure data must label them as predictions, not direct measurements \citep{joo2016}.

\section{Whole-Body Momentum Balance}\label{sec:grf-impulse}\label{grf-whole-body-momentum-balance}

For a golfer-plus-club system of total mass \(m\), before ball impact and after neglecting small aerodynamic loads, the linear momentum balance is

\begin{equation}
\bm{F}_{\mathrm{GRF}}+m\bm{g}=m\ddot{\bm{r}}_{\mathrm{COM}},
\label{eq-grf-com}
\end{equation}

where \(\bm{g}\) points downward. With a vertical axis positive upward and \(g=9.81\ \mathrm{m\,s^{-2}}\),

\[
F_z-mg=m\ddot z_{\mathrm{COM}}, \qquad
\ddot z_{\mathrm{COM}}=\frac{F_z}{m}-g.
\]

For example, \(F_z=1200\) N for a 90 kg system gives \(\ddot z_{\mathrm{COM}}=1200/90-9.81=3.52\ \mathrm{m\,s^{-2}}\) upward. A formula yielding \(23.1\ \mathrm{m\,s^{-2}}\) would have added gravity twice. If \(F_x=200\) N at the same instant, \(\ddot x=2.22\ \mathrm{m\,s^{-2}}\). These are instantaneous accelerations, not velocities or distances. Downward motion can coexist with upward acceleration: braking a descent requires \(F_z>mg\) under this system boundary. Downward acceleration requires \(F_z<mg\), irrespective of whether the system is currently moving up or down.

Integrating Equation~\ref{eq-grf-com} over \([t_1,t_2]\) gives

\[
\int_{t_1}^{t_2}\bm{F}_{\mathrm{GRF}}\,dt
=m\Delta\dot{\bm{r}}_{\mathrm{COM}}-m\bm{g}(t_2-t_1).
\]

This is a strong data-quality check: the measured resultant impulse should agree with the change in whole-system momentum under the same filtering, events, mass, and coordinate convention. Impulse has units N s; mechanical work has units N m. An impulse is not an energy quantity and does not, by itself, specify an ordering of segment-speed peaks. The legs, trunk, arms and club form a coupled system with simultaneous interactions. A proximal-to-distal sequence is an outcome to explain with kinematics, joint powers, inertia and inputs, not a theorem of the whole-body impulse balance. A smaller distal mass alone does not prove greater speed.

The angular-momentum balance about the system COM is

\[
\dot{\bm{H}}_{\mathrm{COM}}
=\sum_i (\bm{r}_{i}-\bm{r}_{\mathrm{COM}})\times\bm{F}_i
+\sum_i\bm{M}_{i,\mathrm{free}}.
\]

Consequently, the same resultant force can have different rotational effects when its line of action, bilateral allocation, or free moment differs. Uniform gravity has zero resultant moment about this COM. If another external wrench is retained, include its moment too. The complete resultant force and moment determine the whole-system momentum rates; they generally do not determine every individual foot wrench or internal joint load.

\section{Contact Work Is a Conditional Statement}\label{sec:grf-work}\label{grf-contact-work-is-a-conditional-statement}

For a rigid contacting body, contact power can be represented at a point \(C\) using its wrench and the body's instantaneous rigid-body velocity evaluated at that same location:

\[
P_c=\bm{F}_c\cdot\bm{v}_c+\bm{M}_c\cdot\bm{\omega}_c.
\]

For an \textbf{ideal rigid, stationary, no-slip contact}, the constrained contact velocity is zero, and the ideal constraint wrench does no work. This useful idealization does not imply zero muscular work. It also does not apply without qualification when the shoe or surface deforms, the support moves, or slip occurs. Ideal rolling without slip on a stationary rigid surface can still have zero contact power; rolling alone is not proof of nonzero work. A COP is an equivalent-wrench location, not necessarily a material point whose velocity can be inserted into a power calculation.

In particular, differentiating the COP coordinates gives the velocity of a moving \textbf{location}, not the contacting body's material velocity there. On a stationary planted rigid foot, the COP can travel while all contact material points remain stationary. For a deformable interface, the general power is the surface integral \(\int_S\bm t\cdot\bm v\,dA\) of traction times local material velocity; one resultant wrench is insufficient without assumptions about the velocity field.

The wrench formula is independent of the reference point when transported consistently. With \(\bm r=\bm r_{OC}\),

\[
\begin{aligned}
\bm v_C&=\bm v_O+\bm\omega\times\bm r,\\
\bm M_C&=\bm M_O-\bm r\times\bm F,\\
\bm F\cdot\bm v_C+\bm M_C\cdot\bm\omega
&=\bm F\cdot\bm v_O+\bm M_O\cdot\bm\omega.
\end{aligned}
\]

The cancellation follows from the scalar triple product. Mixing a moment about one origin with the velocity at another violates this identity.

Muscles and tendons may perform positive, negative, and elastic work while an ideal external contact wrench has zero instantaneous power at its constrained point. GRF can still redirect momentum and alter how internally supplied energy is distributed among segments.

Multiplying the COM balance by \(\dot{\bm r}_{\mathrm{COM}}\) yields

\[
\frac{d}{dt}\left(\tfrac12m\|\dot{\bm r}_{\mathrm{COM}}\|^2\right)
=(\bm F_{\mathrm{GRF}}+m\bm g)\cdot\dot{\bm r}_{\mathrm{COM}}.
\]

This is a translational COM kinetic-energy identity. The term \(\bm F_{\mathrm{GRF}}\cdot\dot{\bm r}_{\mathrm{COM}}\) is not generally the actual contact power. Total mechanical energy also includes segmental motion relative to the COM, gravitational potential, and retained elastic storage. Muscle work and release of previously stored energy can increase kinetic energy while a stationary ideal contact does no work. Deformation may store or dissipate energy, and a moving support may exchange work with the system.

\section{Constrained-Reaction Mechanics}\label{sec:grf-drift}\label{grf-constrained-reaction-mechanics}

At one state, write the multibody equations as

\begin{equation}
\begin{aligned}
M(q)\ddot q&+h_0(q)+h_v(q,\dot q)\\
&=B(q)u+Q_{\mathrm{ext}}+J(q)^\mathsf{T}\lambda.
\end{aligned}
\label{eq-grf-multibody}
\end{equation}

with acceleration constraint

\[
J(q)\ddot q+\gamma(q,\dot q)=0.
\]

Use independent local generalized coordinates \(q\) so that \(M\) is positive definite. If a model uses redundant quaternion coordinates, formulate these equations in independent tangent velocities with the corresponding kinematic map rather than blindly inverting a redundant coordinate mass matrix. Here \(h_0\) is velocity independent, \(h_v\) is velocity dependent, \(u\) is the declared controllable generalized input, \(Q_{\mathrm{ext}}\) contains other known applied loads, \(J\) is the active-contact Jacobian, and \(\lambda\) is the reaction in contact coordinates. Define

\[
A=JM^{-1}J^\mathsf{T}.
\]

For the formulas immediately below, contact constraints are stationary and holonomic, \(\Phi(q)=0\), with \(J=\partial\Phi/\partial q\) and \(\gamma=\dot J\dot q\). Define \(h_0=h(q,0)\) and \(h_v=h(q,\dot q)-h(q,0)\), so \(h_v(q,0)=0\). Hold the declared \(Q_{\mathrm{ext}}\) fixed in the pointwise comparisons. The matrix \(B\) and retained external loads must be independent of \(u\) for the stated affine split.

Solving the motion equation for \(\ddot q\), substituting into the acceleration constraint and collecting the \(JM^{-1}J^\mathsf{T}\lambda\) term gives the following result. If \(J\) has full row rank, \(A\) is positive definite; adequate conditioning is also needed for a reliable numerical solution:

\begin{equation}
\begin{aligned}
d&=h_0+h_v-Bu-Q_{\mathrm{ext}},\\
\lambda&=A^{-1}\left[JM^{-1}d-\gamma\right].
\end{aligned}
\label{eq-grf-lambda}
\end{equation}

The reaction is affine in the declared generalized inputs, even though it is not itself an independently commanded input. A reaction is therefore not automatically a drift-only effect. Eliminating the constraint reactions changes both the drift and input terms of the constrained state equation. It can be separated as

\[
\lambda=\lambda_0+\lambda_v+\lambda_u+\lambda_e,
\]

with

\[
\begin{aligned}
\lambda_0&=A^{-1}JM^{-1}h_0,\\
\lambda_v&=A^{-1}(JM^{-1}h_v-\gamma),\\
\lambda_u&=-A^{-1}JM^{-1}Bu,\\
\lambda_e&=-A^{-1}JM^{-1}Q_{\mathrm{ext}}.
\end{aligned}
\]

These components are exact only for the declared model and state. Errors in kinematics, inertial parameters, contact assignment, external loads, and filtering enter the residual. Dynamic inconsistency is therefore a primary falsification outcome, not a nuisance to hide \citep{sturdy2022, werling2023}.

If constraint rows are redundant, do not use the ordinary inverse of \(A\). Determine the independent constraints and the admissible reaction set. A minimum-norm pseudoinverse can select one solution but cannot establish a unique physical distribution without additional measurements or a justified contact/compliance model.

\subsection{A Coupled Example With Zero Contact Work}\label{sec:grf-remote-torques}\label{grf-a-coupled-example-with-zero-contact-work}

Consider a deliberately simple two-coordinate translational model, with both coordinates measured in meters, \(M\) in kilograms and generalized forces in newtons:

\[
M=\begin{bmatrix}2&1\\1&2\end{bmatrix},\qquad
J=\begin{bmatrix}1&0\end{bmatrix},\qquad
B=\begin{bmatrix}0\\1\end{bmatrix}.
\]

Let \(h_0=h_v=Q_{\mathrm{ext}}=\gamma=0\), and constrain \(q_1=0\). The direct equations are

\[
2\ddot q_1+\ddot q_2=\lambda,\qquad
\ddot q_1+2\ddot q_2=u.
\]

Consequently \(\ddot q_1=0\), \(\ddot q_2=u/2\) and \(\lambda=u/2\). The Schur formula gives the same result because \(A=2/3\) and \(JM^{-1}B=-1/3\). The contact reaction depends on input, yet its power \(\lambda\dot q_1\) is zero. This is an algebraic counterexample, not a calibration of shoulder torque into foot force. A human estimate requires the whole body's inertial coupling, joint geometry, state and active contacts. Planted feet do not in themselves prohibit torso rotation.

\subsection{Pointwise Ground-Reaction ZTCF}\label{grf-pointwise-ground-reaction-ztcf}

The pointwise ground-reaction zero-torque counterfactual (ZTCF) sets the declared controllable input to zero at the achieved \((q,\dot q)\) while retaining declared non-control external loads:

\[
\lambda_{\mathrm{ZTCF}}=\lambda_0+\lambda_v+\lambda_e.
\]

It answers what reaction the model requires at this state with zero declared control. It is not a no-muscle experiment and does not describe what the body would do after controls were removed. The retained state may contain activation or elastic energy from prior action; zero generalized torque and zero muscle excitation are different interventions.

\subsection{Pointwise Ground-Reaction ZVCF}\label{grf-pointwise-ground-reaction-zvcf}

The pointwise ground-reaction zero-velocity counterfactual (ZVCF) evaluates the same configuration and declared inputs after velocity is zeroed:

\[
\lambda_{\mathrm{ZVCF}}=\lambda_0+\lambda_u+\lambda_e.
\]

ZTCF and ZVCF are not forward simulations. They overlap in \(\lambda_0+\lambda_e\) and must not be added to reconstruct the total. The additive attribution is configuration plus velocity plus control plus other external load.

\subsection{Feasibility and Moving Constraints}\label{grf-feasibility-and-moving-constraints}

An equality-constrained calculation is not sufficient to establish a physical contact. A unilateral contact requires a nonnegative compressive normal reaction; sticking additionally requires a tangential force inside the declared friction cone, for example \(\|\bm F_t\|\leq\mu F_n\). A finite support patch also bounds admissible moments and COP. Test these conditions for the full model and for each recomputation. A negative normal reaction means that the retained mode would require adhesion, not that the ground actually pulls the shoe downward. Lift-off or slip requires a different mode.

For the toy model above, interpreting \(q_1\geq0\) as a normal gap makes a negative \(u\) incompatible with the assumed closed, non-adhesive contact because \(\lambda=u/2<0\). Analogous failures can occur when zeroing an input in a loaded model. For example, retain \(Q_{\mathrm{ext}}=(1,0)^\mathsf{T}\) N in the toy model. Then \(\lambda=u/2-1\) N: \(u=4\) N gives an admissible 1 N reaction, while zero input would require \(-1\) N. Releasing the contact instead gives \(\ddot q_1=2/3\ \mathrm{m\,s^{-2}}\), consistent with separation. Report an infeasible fixed-mode diagnostic as such; do not silently relabel it a realizable swing.

For a moving support \(\Phi(q,t)=0\), the velocity constraint is \(J\dot q+\Phi_t=0\), and the acceleration bias contains prescribed-motion terms that may survive at zero velocity. Zeroing all velocities may violate the velocity constraint itself. Reevaluate the complete equations, specify which support motion and external loads are retained, and label any velocity projection as a different intervention. The compact ZVCF expression above is valid only under its stationary-constraint and held-load assumptions. Even for a stationary support, a drag load reevaluated at zero velocity changes \(\lambda_e\) unless the protocol explicitly freezes that load.

\section{What a GRF Residual Cannot Establish}\label{grf-what-a-grf-residual-cannot-establish}

Measured GRF minus modeled ZTCF is a residual containing input-induced reaction, model error, measurement error, and omitted external loads. It cannot uniquely identify muscle torques. Even with a perfect resultant wrench, the model cannot uniquely determine the bilateral allocation when the contact system is redundant. A pseudoinverse supplies a numerical allocation, not new physical information.

With known kinematics, inertia, and individual external wrenches, inverse dynamics estimates the \textbf{net generalized joint load} required by the model. That load can combine muscle, tendon, passive tissue and other modeled effects. It is not a measurement of each muscle force, activation, effort, or injury risk. Differentiation and filtering errors matter, as do unmeasured club loads and contact allocation. Multiple muscle-force combinations can produce the same net joint moment.

\section{Interpreting Common Golf Measurements}\label{sec:grf-practical}\label{grf-interpreting-common-golf-measurements}

\subsection{Vertical Force}\label{grf-vertical-force}

A vertical peak may combine configuration support, velocity-dependent whole- body acceleration, and input-induced reaction. Peak magnitude alone does not identify which term dominates. ZTCF and ZVCF provide model-based tests of which features persist after declared terms are removed.

\subsection{Horizontal Shear}\label{grf-horizontal-shear}

Anterior-posterior and medial-lateral components are particularly sensitive to axis definitions, stance orientation, foot contact, and filtering. Report the laboratory-to-golfer coordinate transform and avoid changing signs implicitly for left-handed participants.

\subsection{Center of Pressure}\label{grf-center-of-pressure}

COP locates an equivalent force together with the residual free couple. It is not the center of mass, not ``weight transfer,'' and not energy transfer. Ball and Best observed more than one recurring COP strategy across skill levels \citep{ballbest2007}. The literature therefore does not justify one universal COP or GRF waveform as the correct pattern.

\subsection{Bilateral Forces and Free Moment}\label{grf-bilateral-forces-and-free-moment}

Two force plates provide more information than a combined resultant, but bilateral decomposition remains sensitive to each foot's frame, crossover, partial contact, and plate assignment. The vertical free moment represents a distributed torsional contact effect. It should not be equated with joint torque or with the two-hand force couple in a club model merely because the signs look similar.

\subsection{Association With Clubhead Speed}\label{grf-association-with-clubhead-speed}

Golf studies have associated selected lead- and trail-foot force/moment features with clubhead speed \citep{han2019}, and shoe-ground torque differs across some skill and footwear conditions \citep{worsfold2008}. These associations do not establish that maximizing a peak will increase speed for an individual. Rachnavy and colleagues reported improved clubhead-speed association when foot-ground variables were combined with trunk and transfer variables \citep{frontiers2026footground}, but their cross-sectional mediation analysis is not a controlled causal intervention. The systematic review documents substantial methodological heterogeneity \citep{watson2026grf}.

For the Rachnavy study, the initial regression block includes skill group as well as COP factors; its explained variance must not be credited to foot variables alone. Its constructed transfer factor is not automatically a dimensionless mechanical efficiency. Establish numerator, denominator, system boundary and power/impulse definitions before interpreting such a quantity as energy conversion. Neither that association nor a force peak establishes the benefit of a particular coaching intervention.

Coaching descriptions such as sway, hanging back and early extension require operational kinematic definitions and synchronized observation. A GRF or COP trace alone does not uniquely identify one of these patterns or prove that it caused a poor strike. Changes should be evaluated against the golfer's task outcome, repeatability and relevant constraints, rather than inferred from agreement with a single reference waveform.

\section{Falsifying a Drift-Based GRF Model}\label{grf-falsifying-a-drift-based-grf-model}

A defensible human study needs synchronized bilateral six-axis force plates, whole-body and club kinematics, participant-specific or uncertainty-bounded segment inertial parameters, and declared filtering and event rules. Model development and threshold selection should be separated from evaluation on held-out participants.

At minimum, report for every force and moment component:

\begin{itemize}
\item
  bias, RMSE, and an explicitly normalized RMSE;
\item
  waveform \(R^2\) with amplitude error reported alongside it;
\item
  vector impulse error over preregistered phases;
\item
  peak magnitude and timing error;
\item
  COP error only above a declared vertical-force threshold;
\item
  free-moment error for three-dimensional contact models;
\item
  residual pelvis forces and moments as dynamic consistency diagnostics; and
\item
  sensitivity to inertial parameters, filtering, coordinate transforms, and contact allocation.
\end{itemize}

First specify the tested hypothesis: for example, a declared zero-input recomputation predicts a particular phase's vertical impulse within a preregistered tolerance. Reject or narrow it when the full model fails its validation, the zero-input prediction fails held-out tests, or the attributed share changes qualitatively under plausible assumptions. Do not convert one failed implementation into a disproof of every possible drift model.

Include a center-of-mass acceleration baseline. For the total resultant force, the COM equation is already an exact balance under its assumptions, not an approximation that added model complexity must beat. A richer model should justify additional claims about individual foot wrenches, internal loads or interventions on appropriately observed outcomes. Compare shared outputs with the same available inputs, filtering, events and error metrics.

Declare whether force measurements were used to fit kinematics, inertial parameters or residual corrections. In-sample agreement after such fitting is a consistency check, not an independent prediction of the fitted force. AddBiomechanics uses measured ground forces in its dynamics optimization \citep{werling2023}; RRA parameter optimization reduces manual tuning of dynamic consistency \citep{sturdy2022}. These methods are useful, but neither method's success alone validates a golf-specific drift attribution. Separate calibration from evaluation, and keep participant-specific calibration data distinct from held-out evaluation trials. In a predictive comparison, infer inputs without using the evaluation wrench that the model is supposed to predict.

\section{Key Takeaways}\label{grf-key-takeaways}

\begin{itemize}
\item
  GRF is an external contact wrench and a constraint reaction, not a direct motor command.
\item
  Whole-body resultant GRF is strongly constrained by center-of-mass acceleration, but bilateral allocation, COP, and free moment require additional spatial contact information.
\item
  Ideal stationary no-slip contact does zero work at the constrained point; this conditional statement does not eliminate muscular, tendon, shoe, or surface work.
\item
  Pointwise reaction ZTCF and ZVCF are overlapping diagnostics, not additive causes and not forward simulations.
\item
  A measured-minus-ZTCF residual combines control-induced reaction with model and measurement error; it does not reveal unique muscle torques.
\item
  Golfers exhibit multiple viable force and COP patterns. Evaluate individual task outcomes and uncertainty, and test claims against evidence.
\end{itemize}

\section{Chapter Exercises}\label{grf-chapter-exercises}

\begin{enumerate}
\def\labelenumi{\arabic{enumi}.}
\item
  A 95 kg golfer-plus-club system stands still, with only gravity and ground forces external. Compute vertical GRF with upward positive.
\item
  For the same system, \(F_z=1350\) N during the downswing. Compute vertical COM acceleration. Can it be moving downward at that instant? Also compute the velocity change if this force remains constant for 0.10 s.
\item
  Explain when an ideal contact wrench does zero work and list three reasons a real shoe-ground interface may violate the idealization.
\item
  Derive Equation~\ref{eq-grf-lambda} from Equation~\ref{eq-grf-multibody} and the acceleration constraint.
\item
  Show algebraically why ZTCF plus ZVCF double-counts configuration and retained external-load reaction. State the assumptions needed for the compact ZVCF.
\item
  Give two different bilateral foot-force allocations with the same combined resultant and explain what additional measurements distinguish them.
\item
  Design a held-out test that compares a COM-only GRF baseline, ZTCF, and the full constrained model without splitting samples from the same participant across training and test sets.
\item
  Recover COP and free moment from the offset-surface numerical example. Explain why differentiating that COP path does not generally give contact power, and why inverse joint moments do not uniquely determine muscle force.
\end{enumerate}

\section{Worked Solutions}\label{grf-worked-solutions}

\begin{enumerate}
\def\labelenumi{\arabic{enumi}.}
\item
  Static vertical balance gives \(F_z=mg=95(9.81)=931.95\) N. If the 95 kg excluded a held club, its hand load would have to be included instead of silently applying the golfer-plus-club equation to golfer mass alone.
\item
  With upward positive,

  \[
  \begin{aligned}
  \ddot z&=1350/95-9.81\\
  &=4.4005\ \mathrm{m\,s^{-2}},\\
  \Delta\dot z&=(1350-931.95)(0.10)/95\\
  &=0.4401\ \mathrm{m\,s^{-1}}.
  \end{aligned}
  \]

  A negative initial vertical velocity is compatible with that positive acceleration: the descent is being braked. Acceleration alone does not reveal the current direction of travel or displacement.
\item
  For an ideal rigid, stationary, no-slip contact, the reaction pairs with zero constrained velocity, so its power is zero. A moving support, material deformation with nonzero local traction power, or tangential slip can invalidate the idealization. Pure rolling without slip need not do so. Muscle or elastic work elsewhere can remain nonzero. A particle constrained to a stationary frictionless curved track provides a simple example: the normal force redirects velocity while remaining perpendicular to it.
\item
  Write \(h=h_0+h_v\). Substitution gives

  \[
  \begin{aligned}
  0&=JM^{-1}(Bu+Q_{\mathrm{ext}}+J^\mathsf{T}\lambda-h)+\gamma,\\
  A\lambda&=JM^{-1}(h-Bu-Q_{\mathrm{ext}})-\gamma.
  \end{aligned}
  \]

  Invert \(A\) only for independent constraint rows. In the two-coordinate example, direct solution and this formula both give \(\lambda=u/2\). If the input channel has more independent components than the measured reaction, the reaction-to-input inverse cannot generally be unique.
\item
  Under the stated stationary-contact, zero-velocity-bias and held-load assumptions,

  \[
  \begin{aligned}
  \lambda_{\mathrm{ZTCF}}+\lambda_{\mathrm{ZVCF}}
  &=2\lambda_0+\lambda_v+\lambda_u+2\lambda_e\\
  &=\lambda+\lambda_0+\lambda_e.
  \end{aligned}
  \]

  They double-count the retained baseline. Moving supports or reevaluated velocity-dependent loads require recomputing the relevant terms, and any retained unilateral/sticking mode must pass its feasibility conditions.
\item
  Place point contacts at \((-a,0,0)\) and \((a,0,0)\), with total vertical force \(F>0\). Allocations \((F/2,F/2)\) and \((3F/4,F/4)\) have the same resultant \((0,0,F)\) but different moments about the midpoint: \(M_y=0\) and \(aF/2\). A measured total moment or separate plates distinguishes these cases. Even a complete total wrench need not fix all individual loads: add equal and opposite horizontal forces directed along the line joining the contacts. Their total force and moment are zero. Separate shear measurements can distinguish this self-equilibrated allocation, provided friction permits it.
\item
  Split by participant before choosing model structure, thresholds or regression parameters. Define any allowed personal calibration separately. Obtain synchronized kinematics and plate data, but reserve the evaluation plate wrench as an outcome. Use kinematics and declared inertial parameters for the COM resultant baseline; use independently specified or estimated inputs for the full constrained prediction and zero the declared channel for ZTCF. Report errors, uncertainty and contact feasibility on matching outcomes. Individual-foot moments provide additional tests beyond the COM resultant. If inverse dynamics used the same force trace to infer the inputs, label the reconstructed full-model agreement as an identity or consistency check and obtain separate evidence for predictive claims.
\item
  With \(z_s=0.04\) m,

  \[
  \begin{aligned}
  x_C&=[0.04(100)-(-76)]/800\\
  &=0.10\ \mathrm m,\\
  y_C&=[-38+0.04(-50)]/800\\
  &=-0.05\ \mathrm m,\\
  x_CF_y-y_CF_x&=0.10(-50)-(-0.05)(100)\\
  &=0,\\
  M_{C,z}&=8-0=8\ \mathrm{N\,m}.
  \end{aligned}
  \]

  For power use the rigid body's velocity field at \(C\), not the rate at which the equivalent COP location changes. For deformable contact, use local traction power or a justified reduction. A net joint moment similarly summarizes combined actions: antagonistic muscle forces can change together while leaving their net moment unchanged, so inverse dynamics alone cannot identify individual muscle forces.
\end{enumerate}
